\section{Discussion}
\label{sec:discussion}

\subsection{What the arch period buys}

The single most useful property of scattered light for an analysis restricted
to public data is that \cref{eq:fdrive} makes the driving frequency a
\emph{strain-only} measurement. No auxiliary channel is required to say that a
particular trigger was caused by motion at \SI{0.15}{\hertz} rather than at
\SI{2}{\hertz}; the arch period says so directly.

That is what allows the O4a half of this study to exist at all. The O4a
auxiliary release contains nothing that witnesses motion
(\cref{sec:data:aux:o4a}), so an approach that depended on auxiliary data
would simply stop at the end of O3. Reading the mechanism off the morphology
does not.

The inference is only as good as the relation it rests on, which is why
\cref{sec:results:aux} matters beyond its own numbers. Establishing in O3 ---
where both the labels and the witnesses exist --- that the arch period does
track ground motion, and that the predictor $f_n = 2n|v|/\lambda$ quantitatively
explains individual arch trains, is what licenses reading an O4a arch period
as a statement about the instrument's environment. Without the O3 auxiliary
work, the O4a arch-period distribution would be a description of pictures
rather than a measurement of physics.

\subsection{Relation to previous work}

\citet{Soni2021Scattering} established the two-population picture in O3 using
internal data, and identified relative motion between the main and reaction
chains as the dominant slow-scattering path. \citet{Nandi2026} report on
scattered light at Livingston during O4, again finding two populations: one
traced to microseismic ground motion and mitigated by installing baffles, the
other to vertical ground motion coupling through a vacuum chamber and
addressed by adding seismic isolation.

This work is not a substitute for either. It has access to a fraction of the
channels, none of the engineering record, and no ability to intervene in the
instrument. What it adds is a different kind of statement: a reproducible,
externally checkable measurement of how the scattered-light population changed
between the two runs, derived entirely from data anyone can download, with
every intermediate step validated against signals of known parameters. Where
its conclusions agree with the internal analyses, that agreement is
independent evidence; where they disagree, the disagreement is informative
about what public data can and cannot support.

The comparison with \citet{Nandi2026} is the most direct test available. If
the arch-period distribution measured here from O4a strain alone reproduces
the two-population structure they identify from internal channels, the
strain-only method is validated on real data as well as on injections.

\subsection{Limitations}

\paragraph{The O4a class definition is inherited from O3.} Every O4a number
rests on a classifier trained on O3 labels, so it presumes that the
morphological boundary of ``scattered light'' did not move between runs. That
presumption is not testable without O4a labels. The cross-validated efficiency
and false-positive rate on held-out O3 folds
(\cref{sec:results:o4a}) are the quantitative statement of how well the
transfer is expected to hold \emph{within} O3; they are an upper bound on how
well it holds across runs, not an estimate of it. A genuinely novel O4a glitch
morphology resembling scattered light would be absorbed into the sample
without comment.

\paragraph{The trigger generator is not Omicron.} The excess-power finder of
\cref{sec:methods:search} is a stand-in. Its efficiency curve differs from
Omicron's, in ways not characterised here. This is why the comparison is run
like-for-like, with the same finder on both runs: the absolute rates are
specific to this pipeline, and only the ratio between runs is meant to be
compared with anything external.

\paragraph{Coverage.} The O4a search covers a few per cent of the run. The
sample is representative in the statistical sense --- windows are drawn
uniformly in livetime and stratified across the run --- but it cannot resolve
episodic behaviour, and a fortnight of anomalous scattering would appear only
as a handful of windows.

\paragraph{One free parameter in the predictor.} The counts-to-metres
calibration is fitted, not derived, because GWOSC publishes auxiliary channels
uncalibrated. The fit is honest --- trained and validated on disjoint sets,
with the constraint width reported --- but a genuinely calibrated witness
would turn the agreement statistic from a goodness-of-fit into a real test of
\cref{eq:fn}.

\paragraph{Cross-validation is not grouped in time.} Triggers minutes apart
are not independent, so the quoted AUC is mildly optimistic. Time-grouped
folds would be the correct treatment and are the obvious next refinement.

\paragraph{\Hone{} has no public motion witness.} The per-trigger predictor is
a Livingston-only result. Whether the same calibration and the same explained
fraction would be found at Hanford is untested, and untestable from public
data.

\subsection{What would strengthen this}

In rough order of value per unit effort:

\begin{enumerate}
  \item \textbf{An O4a trigger catalogue.} Even an unlabelled Omicron trigger
        list would remove the largest systematic here, by decoupling the
        trigger generator from this pipeline.
  \item \textbf{Time-grouped cross-validation}, replacing the stratified folds
        with folds split on contiguous blocks of the run.
  \item \textbf{Full O4a coverage}, which is a matter of compute rather than
        method; the pipeline is resumable and would scale to it unchanged.
  \item \textbf{Direct comparison against \citet{Nandi2026}} on overlapping
        times, which would test the strain-only classification against an
        independent, internally-derived identification.
  \item \textbf{An O4a auxiliary release containing a motion witness}, which
        would make the O4a half of the study symmetric with the O3 half. The
        test suite is written to detect this if it happens.
\end{enumerate}
